\documentclass[11pt,a4paper]{article}
\usepackage[margin=25mm]{geometry}
\usepackage{amsmath,amssymb,amsthm,booktabs}
\newtheorem{lemma}{Lemma}\newtheorem{theorem}{Theorem}
\title{Three involutions and two counts involving order eight}
\author{Human Compression project}\date{}
\begin{document}\maketitle
Let $G=O'N$, and put $g=|G|=460815505920$. For a subgroup $T$, write
$x_T(p,q)$ for the number of ordered pairs $(r,s)$ satisfying
$|r|=p$, $|s|=q$ and $|rs|=2$. Write $e_T(p,q)$ for the number
which generate $T$. We prove the two fixed Connor--Leemans Table~1 entries
\[
 \frac{e_G(3,8)}{2g}=10,\qquad \frac{e_G(6,8)}{2g}=1122.       \tag{1}
\]
The factor $2g$ is the order of $\operatorname{Aut}(G)$; its action on
generating ordered pairs is free. There is no quotient by reversal in (1).
The complete subgroup roster, primary ATLAS data and acquired character
evaluations are borrowed premises. Every row number below refers to that
fixed roster. We reuse the accepted permutation models and containment
arguments, and give the new centralizer and covering argument explicitly.

\section*{The centralizer of an element of order eight}
Let $\mathcal A$ be the conjugacy class of maximal subgroups in row 7.
Its members are exactly the involution centralizers. They are in bijection
with the involutions: a maximal involution centralizer is self-normalizing,
and the single involution class and row 7 have the same class length.

\begin{lemma}\label{three}
For $|s|=8$, the centralizer $C_G(s)$ has exactly three involutions.
Together with the identity they form a group $V_s\cong C_2^2$ containing
$s^4$. The subgroup $V_s$ belongs to row 578, and $|C_G(V_s)|=256$.
\end{lemma}
\begin{proof}
The two element classes of order eight have centralizer order 32.
Each of the two cyclic subgroup classes of order eight, rows 559 and 560,
has a cyclic overgroup of order sixteen, as the complete roster shows.
Choose $a$ of order sixteen with $a^2=s$. Its cyclic group has index two
in $C_G(s)$. If $b$ lies outside it, conjugation by $b$ sends $a$ to $a$
or $a^9$, since it fixes $a^2$. Write $b^2=a^k$. In the first case odd
$k$ would give an element of order 32, absent from the acquired element
list. In the second case, conjugating this equality by $b$ gives
$8k\equiv0\pmod {16}$. Thus $k$ is even in either case. Replacing $b$
by $a^jb$ changes the exponent of its square to $2j+k$ or $10j+k$.
Either congruence can be made zero modulo sixteen. With this choice,
the involutions are precisely $a^8,b,a^8b$, proving the first assertion.

There are exactly two classes of elementary abelian subgroups of order
four. Their normalizers have orders 864 and 1536, for rows 577 and 578
respectively. The only subgroup class of order 864 is row 89. Its
identification as $(C_3^2:C_4)\times S_4$ was proved in the accepted
order twelve component. This group has no element of order eight.
Since $s$ centralizes $V_s$, the normalizer of $V_s$ contains an element
of order eight, so $V_s$ must belong to row 578.

For completeness we identify its full centralizer in the accepted
permutation model $E=4^3.L_3(2)$ of order 10752. In its normal subgroup
$C_4^3$, choose two distinct involutions and let $V$ be the group they
generate. Direct permutation multiplication gives
\[
 |N_E(V)|=1536,\qquad |C_E(V)|=256,
 \qquad N_E(V)^{V\setminus\{1\}}=S_3.                      \tag{2}
\]
The reader checks every element of the retained model and the six
permutations of the three nonidentity elements, each with 256 preimages.
The normalizer order forces $V$ into row 578; the complete normalizer
order in that row then gives $N_E(V)=N_G(V)$. Consequently
$C_E(V)=C_G(V)$. Conjugacy proves the final assertion for every $V_s$.
\end{proof}

The same calculation gives a useful formula beyond the two entries
considered here. Put $C=C_G(V)$ for a row 578 subgroup $V$, and let
$U_{\mathcal A}(p,8)$ count pairs lying in at least one member of
$\mathcal A$. Then
\[
 U_{\mathcal A}(p,8)=[G:A]x_A(p,8)-2[G:N_G(V)]x_C(p,8).       \tag{3}
\]
Indeed the involutions centralizing a pair form the intersection of
$V_s$ with $C_G(r)$, so there are zero, one or three of them. In the last
case the associated group is uniquely $V_s$. Each pair counted three
times in the first term is counted once in the second, with coefficient
two. The literal centralizer has order spectrum
$1,55,104,96$ at orders $1,2,4,8$, and
\[
 x_C(2,8)=1536,\quad x_C(4,8)=3072,\quad x_C(8,8)=672;
\]
all other counts vanish. In particular a pair of type $(3,8)$ or $(6,8)$
is contained in at most one member of $\mathcal A$.

\section*{Two retained finite groups}
\begin{lemma}\label{finite}
Every pair with an entry of order eight and involutory product in
$E=4^3.L_3(2)$ has a common centralizing involution in $E$.
In $J=M_{11}$, every pair of type $(3,8)$ has a common centralizing
involution. For type $(6,8)$, exactly 7920 pairs have such an involution;
the other 47520 pairs generate $J$.
\end{lemma}
\begin{proof}
We supply finite proofs using the previously accepted complete
permutation models of degrees 112 and 11, respectively. The reader
reconstructs each model from its generators and checks its order,
10752 or 7920. It partitions all elements of order eight into conjugacy
classes by literal conjugation. Both models have two such classes;
their sizes are 336 in $E$ and 990 in $J$.
For one representative $s$ in each class, every possible partner $r$
is tested by direct multiplication. All involutions commuting with
$s$ are explicitly listed; the common involutions are obtained by
testing commutation with $r$. Conjugacy transports the result to the
whole class. The resulting totals needed here are
\[
\begin{array}{c|rr|rr}
 &\multicolumn{2}{c|}{x_T(p,8)}&
   \multicolumn{2}{c}{\text{with a common involution}}\\
 T&p=3&p=6&p=3&p=6\\\hline
 E&21504&21504&21504&21504\\
 J&7920&55440&7920&7920.
\end{array}                                                   \tag{4}
\]
For each uncovered representative pair of type $(6,8)$ in $J$, literal
closure under its two generators produces the entire retained set of
7920 permutations. This proves generation, rather than inferring it
from an order label. The same test was retained for the other uncovered
orders in the diagnostic, but those results are not needed for (1).
For $E$, the complete partner lists for both order eight classes have
no uncovered pair of any order, proving the first assertion.
\end{proof}

\section*{The remaining overgroups}
Write $M=L_3(7):2$. Its two conjugacy classes are rows 2 and 3.
The accepted block and orthogonal centralizer calculations give
\[
 H=C_M(t_{\rm in}),\quad |H|=1344,\qquad
 K=C_M(t_{\rm out})\cong C_2\times\mathrm{PGL}_2(7),\quad |K|=672.
\]
There are 2793 inner and 5586 outer involutions in $M$.
We shall use the already proved unique-parent incidences for rows
4, 5, 9, 10, 13, 14, 45, 46, 172 and 174, and the identification of
$K$ with rows 110 and 111.

\begin{lemma}\label{inside}
For $p=3$ or $6$, if a pair of type $(p,8)$ lies in a projective maximal
$M$ and in an involution centralizer $C_G(t)$, then $t\in M$.
\end{lemma}
\begin{proof}
Put $L=\langle r,s\rangle$. The complete common ancestor join of the
cyclic $p$ and cyclic eight classes, restricted to descendants of
rows 7 and 2 or 3, gives the same list for both values of $p$:
\[
65,66,110,111,112,113,115,116,117,118,171,173,175,176,
297,298,382,383,384,385.                                      \tag{5}
\]
All but rows 110, 111, 171 and 173 have a block overgroup of row 65
or 66. The accepted literal block calculation proves that every
pair with an entry $s$ of order eight commutes with $s^4$.
Thus, for these classes, $s^4\in L$ is a common centralizing
involution. The uniqueness following (3) gives $t=s^4\in M$.

If $L$ is in row 110 or 111, the central involution of its $C_2$
factor gives $t\in L$. For row 171 or 173, the accepted direct
incidences give a unique actual parent $K$ of row 110 or 111, and
a unique simple parent $S$ of row 4 or 5. Here $L$ has index two
in $K$, and $|N_G(L)|=672$, so $K=N_G(L)$. It normalizes the unique
$S$. Its normalizer $N_G(S)$ is the unique projective maximal
parent of $S$. Every path from $L$ to a projective maximal starts
through $K$ or $S$, so that maximal contains $K$. The central
involution of $K$ centralizes the pair and must be $t$. This proves
the assertion in all cases of (5).
\end{proof}

\begin{lemma}\label{outside}
Outside $\mathcal A$, a nongenerating pair of either type lies in
exactly one projective maximal, generates a row 15 or 16 subgroup
$L_2(31)$, or generates a row 23 or 24 subgroup $M_{11}$.
These alternatives are disjoint. The $L_2(31)$ alternative occurs
only for $(3,8)$; the $M_{11}$ alternative occurs only for $(6,8)$.
\end{lemma}
\begin{proof}
The exact character evaluations show that the only maximal classes
with nonzero counts are rows 2, 3, 7, 15, 16, 21, 23 and 24 for
$(3,8)$, and rows 2, 3, 7, 21, 23 and 24 for $(6,8)$.
Lemma~\ref{finite} removes row 21 outside $\mathcal A$ and gives
the precise assertion for rows 23 and 24.

For clarity we also check that no additional projective overlap is
hidden in these reductions. Intersect the complete ancestors of the
cyclic $p$ classes with those of rows 559 and 560. Remove $G$ and
all descendants of rows 7, 11, 12, 21, 23 and 24. The remainder is
\[
2,3,4,5,9,10,13,14,45,46,172,174,                             \tag{6}
\]
with rows 15 and 16 added for $p=3$. Removing rows 11 and 12 is
legitimate because their pair counts are zero. Removing a proper
subgroup of an $M_{11}$ is legitimate by Lemma~\ref{finite}; a pair
outside $\mathcal A$ there would have to generate the full $M_{11}$.

The nonmaximal rows in (6) have, successively, unique actual parents
\[
(4,2),(5,3),(9,4),(10,5),(13,9),(14,10),(45,13),(46,14),
(172,4),(174,5),                                             \tag{7}
\]
where each ordered pair means child and parent. For each edge the
complete class-length incidence equals one. Composing these proved
uniqueness assertions gives exactly one projective maximal. No other
maximal class can occur above one of these rows. A pair generating a
listed maximal subgroup cannot also lie in a different maximal.
Finally the cyclic ancestor join for $L_2(31)$ has no proper subgroup
containing both a three-element and an eight-element. Thus all of
its $(3,8)$ pairs generate it. This proves the lemma.
\end{proof}

\begin{theorem}
The two numbers in (1) follow from one common subtraction formula.
\end{theorem}
\begin{proof}
The two projective classes each have 122760 members. By
Lemma~\ref{inside} and the uniqueness of the common involution,
the number of pairs in one projective maximal which are covered by
$\mathcal A$ is
\[
                       2793x_H(p,8)+5586x_K(p,8).             \tag{8}
\]
Write $n_T$ for the length of a subgroup class. By
Lemmas~\ref{finite}--\ref{outside}, the number of nongenerating pairs is
\[
 n_Ax_A(p,8)+\sum_{i=2,3,15,16}n_i x_i(p,8)
 +\sum_{i=23,24}n_i e_J(p,8)
 -2\cdot122760\bigl(2793x_H(p,8)+5586x_K(p,8)\bigr).          \tag{9}
\]
The needed block counts are $672,6048$ and the orthogonal counts are
$1344,4032$, in order $p=3,6$. The finite lemma gives $e_J=0,47520$.
All other entries of (9) are the retained character evaluations.
Substitution gives
\[
\begin{array}{c|rr}
 p&x_G(p,8)&\text{nongenerating pairs}\\\hline
3&24423221813760&15206911695360\\
6&1156186104353280&122116109068800.
\end{array}
\]
Subtracting the second column from the first and dividing by $2g$
gives 10 and 1122. The arithmetic reader also compares these with
the fixed printed cells and retains the thirteen maximal contributions.
\end{proof}

\paragraph{Attribution and limits.}
The ATLAS group identifications, complete subgroup classification,
character tables and the accepted small permutation models are imported.
The new proof consists of Lemma~\ref{three}, the finite covering
certificates, their exact source support and the resulting formula (9).
The suggestion to control overlaps through the involutions of $C_G(s)$
came from a separate project agent; the explicit cyclic-index-two
argument and retained-model applications are supplied here.
This proves two numerical cells. It does not establish any new chirality
claim or complete the fixed original. Formula (3) is available for the
remaining even partner orders, but their other overgroups still require
separate accounting.
\end{document}
